Design in power electronics—encompassing topology selection, parameter tuning, and reliability assurance—often involves high-dimensional, nonlinear optimization problems, further complicated by tight constraints. Traditional design workflows—relying on expert experience, manual iterations, and slow simulations—typically lead to suboptimal convergence and prolonged design cycles. This is particularly evident in topology discovery and component coordination, where conventional methods struggle with efficiency and adaptability.

Recent advancements in artificial intelligence (AI), particularly deep reinforcement learning (DRL), offer promising solutions to these challenges. DRL introduces a data-driven approach capable of autonomously learning design strategies by interacting with its environment. Unlike traditional heuristic-based optimization methods, DRL enables adaptive exploration of complex design spaces, effectively handling sparse rewards and multi-objective trade-offs. This adaptability makes DRL particularly well-suited for optimizing two critical sub-tasks in PEC design: topology exploration and parameter optimization.

\subsection{Topology Exploration}
\label{sec:topo_explore}

In the context of PECs, DRL has shown significant potential in automating the topology design of multiport DC–DC converters. The design process is framed as a Markov Decision Process (MDP), where an agent sequentially connects components to form valid circuits, ensuring both electrical feasibility and optimal performance. For instance, a graph-based framework combining Graph Neural Networks (GNNs) with reinforcement learning was proposed to handle the complexities of multiport DC-DC converters~\cite{liang2022learning}. By representing circuits as graphs, the framework captures intricate dependencies, with a PPO agent efficiently connecting components. This method overcomes combinatorial explosion and reduces time for discovering the first feasible circuit from exponential to near-linear scaling, while reducing switch counts by over 30\% compared to conventional solutions. Moreover, simulations demonstrated stable output voltages with minimal error under duty cycles, showcasing the real-world effectiveness of this approach.

Building on this work, another DRL-based approach~\cite{dong2023topology} integrated PPO into a GNN-based framework, achieving significant improvements in topology generation speed and voltage stress management. This framework reduced the topology generation time for 8-port converters by more than 14 times, achieving 50\% lower voltage stress and 31\% lower current stress, optimizing multi-port DC-DC converters for better performance. These advancements highlight DRL's ability to streamline and accelerate topology exploration, significantly improving design efficiency and performance.

\subsection{Parameter Optimization}
\label{sec:param_opt}

Parameter optimization is a crucial aspect of power electronic converter systems (PECS), particularly in balancing efficiency, power density, thermal performance, and cost. DRL has demonstrated its capability to automate the optimization of key parameters, significantly improving both performance and design efficiency.
In \cite{wang2024case, wang2022efficiency}, a DRL-based framework was proposed to optimize three-level ANPC inverters, addressing challenges such as current ripple, thermal constraints, and BOM cost. The optimization problem was framed as a continuous Markov Decision Process (MDP), where states represented operational specifications (\(U_{\text{DC}}, I\)) and actions controlled the switching frequency (\(f_{\text{sw}}\)). This approach led to significant improvements, achieving 18× faster optimization (5s vs. 90s per configuration) and a 7.73\% reduction in losses at 400V/30A. Experimental validation on a 140 kW prototype confirmed greater than 99.8\% efficiency, with only a 0.14\% deviation from theoretical values. These results highlight DRL's potential to substantially accelerate the optimization process while maintaining high accuracy.

Similarly,~\cite{wang2023efficiency} applied DRL to LLC resonant converters, simultaneously optimizing efficiency and power density under dynamic design constraints. The algorithm's state-action-reward framing allowed for the efficient balancing of switching losses and magnetic volume, achieving \textbf{96.57\% efficiency} at 400V/3~kW with \textbf{31.3~kW/dm³ power density}, outperforming traditional design methods. The optimization process was \textbf{18× faster} than brute-force approaches, and the algorithm demonstrated robust adaptability, handling ±50\% power/voltage variations instantly.

Additionally,~\cite{zhang2024reinforcement} tackled the challenge of voltage instability and sluggish dynamic response in Dual Active Bridge (DAB) converters under abrupt load changes. The proposed DQN and DDQN algorithms optimized phase-shift modulation parameters to minimize overshoot and voltage deviations, achieving \textbf{zero overshoot} and \textbf{±0.5\% voltage deviation} during transients, with a sub-millisecond response to load steps. These results highlight the power of DRL in addressing the dynamic control and parameter optimization challenges of PECs.

For converters like Buck and Boost converters, DRL-based methods such as PPO have shown potential for optimizing parameters like switching frequency, inductance, and capacitance. In~\cite{bui2023deep}, PPO was applied to efficiently handle hardware parameter design, achieving \textbf{97.7\% power efficiency} for Boost converters, \textbf{5× faster computation} (15 hours vs. 72 hours for brute-force), and enabling seamless handling of multi-parameter optimization without model dependency. This work demonstrates DRL's efficiency in solving complex topological optimization problems in PEC systems.

Moreover, advances in magnetic component design also benefit from DRL. For instance,~\cite{tian2022parameter, tian2024optimizing} used DDPG to automate the design of DC inductors under triangular waveforms, achieving sub-second inference and significantly reducing computation time. The algorithm ensured \textbf{97.73\% peak efficiency} with \textbf{0.18\% voltage ripple}, providing a solution to multi-objective optimization complexity and reducing design time by \textbf{48×} compared to brute-force methods.

To address the computational cost and scalability challenges, surrogate-assisted DRL methods, such as those in~\cite{bui2022deep}, integrate DNN-based surrogate models with DRL to accelerate optimization across multiple converter topologies. This method reduced training time by \textbf{36× faster optimization} (30 mins vs. 18 hrs per iteration) while maintaining \textbf{99.26\% power efficiency} for buck converters, demonstrating the feasibility of DRL for high-dimensional design spaces.

Collectively, these advances illustrate DRL's role as a scalable, adaptable tool for parameter optimization, capable of improving the efficiency, speed, and precision of power electronic converter systems. By automating complex design processes and overcoming traditional limitations, DRL offers a significant pathway for advancing the next generation of intelligent PECs.


\subsection{Discussion}
\label{sec:design_conclusion}

In conclusion, DRL has demonstrated its transformative potential in the design and optimization of power electronic converter systems, particularly in the automation of topology exploration and parameter optimization. DRL techniques, such as PPO, DDPG, and SAC, have shown remarkable success in improving efficiency, reducing losses, and accelerating the design process, making them essential tools for modern PECs. The ability of DRL to handle complex, multi-objective optimization tasks allows for faster and more precise designs, offering significant advancements over traditional optimization methods. However, there are several challenges that remain. While DRL methods excel in specific applications, they still face issues such as high computational cost, limited scalability to large systems, and challenges related to sparse rewards during training. Additionally, there is a need for further research into enhancing the robustness and generalizability of these algorithms across diverse power electronic systems. Moreover, integrating DRL with traditional methods in a hybrid approach could help leverage the strengths of both techniques, ensuring the reliability and performance of PEC designs in real-world applications. Future work should focus on addressing these limitations, improving the adaptability of DRL to more complex, large-scale systems, and exploring its integration with domain-specific knowledge to achieve more interpretable and physically grounded solutions. 
